\section{Testing the integer relation without proxy coordinates}
\label{sec:direct-integer-relation}

The first new improvement removes the three rational Lucas proxies
from the shared-point construction.  Two were retained coordinates;
the third was a finite quadratic root.  Their role was to impose
one integer polynomial relation on three elliptic-group labels.
The relation can instead be tested directly on the approximate
logarithmic ratios.  This gives a conditional definition with
$41$ total rational coordinates, including the input, or $43$
with the semilocal anchor.

\subsection{An elementary substitution principle}

\begin{lemma}\label{lem:direct-polynomial-congruence}
Let $L$ be a number field, let $\mathfrak p$ lie over a rational
prime $p$, and normalize its valuation by $v_{\mathfrak p}(p)=1$.
Suppose $z_i\in L$ and $n_i\in\mathbb Z$ satisfy
\[
 v_{\mathfrak p}(z_i)\ge0,
 \qquad v_{\mathfrak p}(z_i-n_i)\ge k
 \quad(1\le i\le r).
\]
For every $F\in\mathbb Z[T_1,\ldots,T_r]$,
\[
 v_{\mathfrak p}\bigl(F(\mathbf z)-F(\mathbf n)\bigr)\ge k.
\]
Consequently a requirement $v_{\mathfrak p}(F(\mathbf z))\ge k$
is equivalent to $p^k\mid F(\mathbf n)$ when $k$ is a positive
integer.  The same congruence assertion holds for rational-coefficient
polynomials whose coefficient denominators are units at $p$; in that
case the final condition is $v_p(F(\mathbf n))\ge k$.
\end{lemma}
\begin{proof}
For each monomial, replace its factors one at a time.  Every term
of the resulting telescoping difference contains some $z_i-n_i$;
all its other factors, including the coefficient, are integral at
$\mathfrak p$.  Thus every term has valuation at least $k$.
The nonarchimedean triangle inequality proves the assertion.
The final equivalence follows because $F(\mathbf n)$ is a rational
integer and the normalized valuation restricts to $v_p$ on it.
\end{proof}

This principle does not make the ratios equal to integer labels.
It transfers a polynomial relation modulo the selected prime powers.
The global height comparison remains the step that forces the
integer residual to vanish exactly.

For a quantitative general version, write
$\|F\|_1$ for the sum of the absolute values of its coefficients.
If $\deg F=d$, $|n_i|\le B$, and $B\ge1$, then
\[
                    |F(\mathbf n)|\le\|F\|_1B^d.
\]
Thus a height argument of the form
\[
 B^2-6<7H|F(\mathbf n)|^{24/k}
\]
forces $F(\mathbf n)=0$ whenever
\[
 k>12d,\qquad
 B>\left((7\|F\|_1^{24/k}+6)H\right)^{1/(2-24d/k)}.
\]
Indeed, if the integer residual were nonzero, the first inequality
and $B,H\ge1$ would imply the reverse strict inequality on $B$.
For several integer relations one can either test them separately,
or use the single relation $\sum_jF_j^2=0$.  The latter doubles
the largest polynomial degree and its precision requirement must
be adjusted accordingly.

\subsection{The replacement in the shared construction}

Retain the nine nonidentity points
\[
              A,S,T,Z,H,U_A,U_S,U_T,U_Z\in2E(K),
\]
where $K=\mathbb Q(\sqrt2)$, and all the point, subtraction,
logarithm and denominator guards of the preferred shared source.
Their root labels are $(n,a,t,b)$ for $(A,S,T,Z)$, respectively.
These are integers used in the proof; they are not rational
coordinates of the polynomial.  Put $\beta=X^2+2$.

Let $L_i,L'_i$, $1\le i\le4$, be the primary and conjugate
alternate logarithmic ratios from the previous construction,
with orders $(16,64,64,64)$.  Delete the rational coordinates
$\xi,\zeta$ and the monic quadratic root $\tau$, together with
the equation involving them.  Define
\[
 F(U,V,W)=V^2-U(U-2)VW+(U-2)^2W^2-1.
\]
Use the following sixteen field tests in every coordinate detector:
\[
\begin{array}{c|c}
 f(U_0),\ldots,f(U_4),\sigma f(V_0),\ldots,\sigma f(V_4)
       &1\quad(10\text{ tests})\\
 L_i-L'_i\quad(1\le i\le4)&16,64,64,64\\
 L_1-\beta&16\\
 F(L_2,L_3,L_4)&64.
\end{array}
\]
Here $U_0=H$, $Q_0=P$, $(Q_1,Q_2,Q_3,Q_4)=(A,S,T,Z)$,
and $V_i=U_i-[4]Q_i$, exactly as before.  There are ten tests
of order one, two of order sixteen, and four of order sixty-four.
Every other part of the preferred source is retained: its fixed
positive square scales, nine reciprocal-trace detectors, normalized
scalar-height block sharing $H$, conic parametrizations, and the
chosen size anchor.

\begin{proposition}[Preservation of the integer input projection]
\label{prop:direct-Lucas-preservation}
Under the explicitly stated arithmetic inputs of the shared
construction, the modified system has input projection exactly
$\mathbb Z$.
\end{proposition}
\begin{proof}
Consider first a rational solution and a prime counted by one
of the nine height radicals.  The detector selects a place
above that prime at which all sixteen tests have their assigned
orders.  The two-anchor argument uses only the first ten tests
and the four differences $L_i-L'_i$.  Hence it still proves
\[
 L_i\text{ is integral},\qquad
 L_i\equiv n_i\pmod{p^{k_i}},
 \qquad(n_1,n_2,n_3,n_4)=(n,a,t,b).
\]
In particular this conclusion does not require any proxy equation.

Lemma~\ref{lem:direct-polynomial-congruence} now gives
\[
 F(L_2,L_3,L_4)\equiv
 \Delta:=t^2-a(a-2)tb+(a-2)^2b^2-1\pmod{p^{64}}.
\]
The new order-$64$ test therefore says exactly that the integer
$\Delta$ is divisible by $p^{64}$.  The radical divisibilities
used in the original height argument are unchanged.  If
$B=\max(|n|,|a|,|t|,|b|)$ and $\Delta\ne0$, that argument
gives
\[
 B^2-6<7H|\Delta|^{3/8},\qquad |\Delta|<5B^4.
\]
The unchanged anchor contradicts these estimates.  Therefore
$\Delta=0$, and the nonzero root label $b$ yields the same
Lucas growth inequality $\log B>|a|-3$.

The normalized scalar block is unchanged and still gives
$h(\beta)\le H\sqrt{|a|}$.  The order-$16$ test $L_1-\beta$
and the corresponding cross comparison give the same divisibility
of the numerator of $\beta-n$.  If $\beta\ne n$, the resulting
bound
\[
                    \log B<2\log(32H)+3h(\beta)
\]
contradicts the anchor and the Lucas growth inequality as before.
Thus $\beta=n\in\mathbb Z$, and $X\in\mathbb Z$.

For completeness, start with an integer input.  Choose the same
integer root labels as in the previous completeness proof, so
$F(a,t,b)=0$.  At every selected place, choose the common index
and points so that $L_i\equiv n_i\pmod{p^{k_i}}$ and the same
holds for the alternate ratios.  The lemma then gives
$v_{\mathfrak p}(F(L_2,L_3,L_4))\ge64$, which is the only
requirement on this new test.  Exact vanishing of $F(L_2,L_3,L_4)$
is not required.

The new test is a fixed polynomial in the old ratios.  Its real
values are bounded on the same compact sets avoiding the old
denominators, and its valuations have finite lower bounds in the
same small-prime neighborhoods.  Consequently the normalization,
prime-pattern, local-square and preparation-zero completeness
arguments apply without alteration.  No new denominator or
exceptional prime is introduced by $F$ itself.  This constructs
all the remaining witnesses for every integral input.
\end{proof}

\subsection{Exact coordinate and finite-algebra counts}

With the finite size exclusion, the retained coordinates are
\[
 \underbrace{1}_{X}+
 \underbrace{18}_{\text{nine halfpoint abscissa pairs}}+
 \underbrace{1}_{\lambda}+
 \underbrace{18}_{\text{nine pole pairs}}+
 \underbrace{2}_{\text{scalar pole pair}}=40.
\]
The finite algebra has nine quartic point generators and ten
quadratic pole generators.  Its rank is at most
\[
                            4^9 2^{10}=2^{28}.
\]
The finite-algebra construction adds one final rational coordinate,
giving $41$ total coordinates and $40$ existential witnesses.
The semilocal anchor adds two retained coordinates and one
quadratic generator, giving rank at most $2^{29}$ and totals
$43$ and $42$, respectively.  The finite correction and conic
choices are unchanged.  The removed Lucas root has no remaining
equation, repeated-root stratum, or denominator guard.

The files \texttt{source-41-direct.json.gz} and
\texttt{source-43-direct-semilocal.json.gz} record the complete
coefficient DAGs and finite parameters.  The copied producer is
\texttt{build\_direct\_source.py}.  The frozen earlier source
files are unchanged.  Final degree and coefficient-height bounds
come from the separately checked finite-algebra compilation of
these new inputs.
